\subsubsection{Hyperbolic Tangent and Other Activations}\label{sec:hyperbolic-tangent-and-other-activations-cnn}

This section is mostly a \textbf{comparison} rather than a new mechanism. Here we will introduce \textbf{tanh} and \textbf{sigmoid} as other valid activation functions, but we will also make clear that, in CNNs, ReLU is the usual practical choice.

\highspace
The \textbf{hyperbolic tangent} function is defined as (\autopageref{sec:tanh-activation-function}):
\begin{equation*}
    \tanh(x) = \frac{2}{1 + e^{-2x}} - 1
\end{equation*}
And its output range is $\left(-1, 1\right)$. So unlike ReLU, which can produce any nonnegative value, $\tanh$ \textbf{squashes} every input into a bounded interval. For example:
\begin{equation*}
    x \to +\infty \quad \Rightarrow \quad \tanh(x) \to 1
\end{equation*}
While:
\begin{equation*}
    x \to -\infty \quad \Rightarrow \quad \tanh(x) \to -1
\end{equation*}
Also unlike ReLU, $\tanh$ is \textbf{differentiable} everywhere, continuous, and bounded.

\highspace
\textbf{Sigmoid} behaves similarly (\autopageref{sec:sigmoid-activation-function}):
\begin{equation*}
    S(x) = \frac{1}{1 + e^{-2x}}
\end{equation*}
But its range is: $(0, 1)$.

\highspace
However, in general sigmoid and tanh are more typical of MLP architectures than CNNs. Let's compare the four activation functions we have discussed so far:

\begin{center}
    \begin{tabular}{@{} l c c c @{}}
        \toprule
        \textbf{Activation} & \textbf{Output range} & \textbf{Negative outputs?} & \textbf{Bounded?} \\
        \midrule
        ReLU        & $\left[0, +\infty\right)$         & No    & No \\
        Leaky ReLU  & $\left(-\infty, +\infty\right)$   & Yes   & No \\
        $\tanh$     & $\left(-1, 1\right)$              & Yes   & Yes \\
        sigmoid     & $\left(0, 1\right)$               & No    & Yes \\
        \bottomrule
    \end{tabular}
\end{center}

\noindent
The practical reason \hl{ReLU became dominant is related to training.} Tanh and sigmoid \textbf{saturate} for large positive or negative inputs: once they are \hl{near their asymptotes, their derivatives become very small.} That can make gradients shrink during backpropagation and contribute to the vanishing-gradient problem (\autopageref{def:vanishing-gradient-problem}).

\highspace
ReLU instead has a much simpler behavior on the positive side:
\begin{equation*}
    \mathrm{ReLU}(x) = x \quad \text{for } x > 0,
\end{equation*}
With derivative 1 there.

\highspace
Tanh and sigmoid are valid nonlinear activations, but \hl{ReLU became the standard choice for CNN hidden layers because it is simpler and generally easier to optimize in deep networks.}

\highspace
Finally, just like ReLU, \hl{tanh and sigmoid are still applied \textbf{element-wise}.} So if the input activation volume is $H \times W \times C$, the output remains $H \times W \times C$. \textbf{Only the values change}.
\begin{itemize}
    \item \textbf{Sigmoid}: works, but saturates and can cause vanishing gradients.
    \item \textbf{Tanh}: generally better than sigmoid around zero because it is zero-centered, but it still saturates.
    \item \textbf{ReLU}: usually preferred because it is simple, fast, and does not saturate for positive inputs.
    \item \textbf{Leaky ReLU}: fixes one important ReLU weakness by allowing a small gradient for $x < 0$, reducing the \textbf{dying ReLU} problem.
\end{itemize}
